% =====================================================================
% Chapter 6: Conclusion and Future Work
% =====================================================================
\chapter{Conclusion and Future Work}
\label{ch:conclusion}

\section{Conclusion}

This thesis developed, implemented, and evaluated a two-part
deep-learning system for jersey number recognition in football images
and broadcast video.

The \textbf{image-based pipeline} combined YOLOv8x person detection,
YOLOv8x-Pose-based torso localization from shoulder and hip keypoints,
a multi-stage image enhancement module, and a multi-variant EasyOCR
engine incorporating dynamic zoom, multiple binarization strategies,
per-digit isolation, projection-based digit splitting, and a
shape-based 0/6/8/9 and 1/7 disambiguation module. A ResNet-18-based
\textbf{legibility classifier} trained on weakly-labeled crops
achieved a best validation accuracy of \textbf{94.55\%}, an ROC-AUC
of \textbf{0.955}, and an average precision of \textbf{0.993}, with
97.87\% accuracy on legible crops and 75.00\% on illegible crops.

The \textbf{video-based pipeline} processed ten broadcast football
clips using a domain-specific \textbf{YOLOv5 football detection model}
combined with the \textbf{ByteTrack} tracker, a temporal role voter for
referee identification, four-variant PARSeq torso processing with an
automatic EasyOCR fallback, and a
\textbf{temporal jersey voting mechanism} over a 60-frame sliding
window. Evaluated against manually annotated ground truth, the system
achieved an average \textbf{tracking accuracy of 95.16\%} (IDF1:
0.692) across 10 videos and a \textbf{GT OCR accuracy of 74.78\%} per
track, rising to \textbf{84.53\%} after applying temporal voting on the
same evaluation rows --- an absolute improvement of 17.4 percentage
points over single-frame predictions. The overall pipeline accuracy
averaged \textbf{79.72\%} across all ten clips.

These results confirm that \textbf{accurate spatial localization},
\textbf{targeted image enhancement}, \textbf{multi-strategy OCR with
shape-based refinement}, and \textbf{temporal aggregation} are all
necessary, complementary components of a practical jersey-number
recognition system, and that no single component alone is sufficient.
Temporal voting, in particular, consistently delivered the largest
single performance gain across all clips, a finding that aligns with
and independently validates conclusions drawn in prior work
\cite{gerke2015soccer,li2018jersey}.

\section{Limitations}

The following limitations of the current system should be acknowledged:

\begin{enumerate}[noitemsep]
    \item \textbf{Limited image evaluation scale:} The image pipeline
    was evaluated on 353 player-image folders rather than the full
    SoccerNet dataset (over 1.2 million images) due to computational
    constraints.
    \item \textbf{Computational cost of detection and OCR:} Running a
    full-resolution YOLOv5 detection pass on every video frame, combined
    with PARSeq/EasyOCR inference on four enhanced variants for every
    detected player in every frame, results in a high total processing
    cost (Section~\ref{sec:results-video}), making the current video
    pipeline unsuitable for real-time deployment without further
    optimization.
    \item \textbf{Weak legibility labels:} The legibility classifier was
    trained on OCR-confidence-derived labels rather than human-verified
    annotations, which may introduce label noise.
    \item \textbf{Digit-range assumption:} The system handles only
    1--2 digit jersey numbers in $[1, 99]$; non-standard markings or
    three-digit numbers are not supported.
    \item \textbf{Fixed torso fraction:} The video pipeline uses a
    fixed-fraction heuristic for torso extraction rather than
    pose-driven keypoints, which may miss digit regions when the player
    is at an oblique angle.
    \item \textbf{Limited clip diversity:} The ten evaluated clips
    share similar broadcast characteristics; generalization to different
    leagues, camera positions, or stadium lighting has not been
    systematically tested.
\end{enumerate}

\section{Future Work}

Building on the findings and limitations identified above, the following
directions are proposed for future investigation:

\begin{enumerate}
    \item \textbf{Purpose-trained digit classifier:} Replace or
    supplement the video pipeline's PARSeq recognizer with a small CNN
    trained directly on cropped jersey-digit images from SoccerNet. This
    would likely improve both accuracy and inference speed by avoiding
    the overhead of a general-purpose scene-text recognizer.

    \item \textbf{Pose-driven torso extraction in the video pipeline:}
    Integrate YOLOv8-Pose into the video pipeline to derive keypoint-based
    torso boxes, replicating the more accurate localization strategy of
    the image pipeline.

    \item \textbf{Evaluation at full dataset scale:} Re-run the image
    pipeline on the complete SoccerNet SN-Jersey dataset ($>$1.2 million
    images) to obtain statistically robust accuracy figures and to
    provide more training data for the legibility classifier.

    \item \textbf{Real-time optimization:} As shown in
    Section~\ref{sec:results-video}, detection now accounts for the
    largest share of total processing time (56.2\%), ahead of OCR
    (42.8\%). Reducing the per-frame detection cost --- for example by
    running detection on a downscaled frame, using a lighter YOLOv5
    variant, or detecting on every second or third frame and relying on
    ByteTrack's motion prediction to bridge the gaps --- combined with
    batching OCR calls across all crops in a frame, is the most direct
    path toward real-time operation.

    \item \textbf{Broader video diversity:} Evaluate the video pipeline
    on clips from multiple leagues, broadcast angles, and stadium
    lighting conditions to assess generalization.

    \item \textbf{Multi-camera tracklet fusion:} Aggregate jersey-number
    evidence across multiple camera angles or over longer tracklets
    by incorporating player re-identification, increasing the evidence
    available to the voting mechanism.

    \item \textbf{End-to-end fine-tuning:} Jointly fine-tune the
    detection backbone and digit recognition head on the jersey-number
    task rather than relying exclusively on off-the-shelf pretrained
    weights.
\end{enumerate}

\section{Final Remarks}

This thesis has demonstrated that combining a domain-specific football
detection model, reliable multi-object tracking, targeted image
enhancement, multi-strategy OCR with shape-based refinement, and
temporal aggregation can produce a meaningfully effective jersey number
recognition system for both static images and broadcast video.
The consistent 17-percentage-point gain from temporal voting
reinforces a key principle that is well-supported in prior literature
and has practical implications for any system aiming to identify players
from broadcast footage at scale.

The complete source code developed for this thesis, covering both the
image-based and video-based pipelines, has been made publicly
available in a source-code repository, the link to which is provided
in the Abstract.
